\textbf{Question and model.} We tested whether machine learning detects stressed markets better than the trend
filter. A two-state Gaussian hidden Markov model (HMM; \citealp{hamilton1989,rabiner1989}) learns a calm and a
turbulent regime from daily S\&P 500 returns and volatility since \res{hmm_data_start}; regime models of this kind are
established in asset allocation \citep{ang2002,nystrup2018}. To avoid look-ahead, the model is re-estimated every
December on past data only, and each month uses the \emph{filtered} probability of the turbulent state given data up
to that day. Smoothed or Viterbi state paths, which use future data and make many published HMM backtests look better
than they are, are never used for decisions. In risk-off months the growth assets (SPY, EFA, VNQ, DBC) are halved.
A simple volatility rule serves as non-ML baseline. Before the first run we fixed all thresholds and three promotion
gates: Sharpe ratio at least that of the base strategy, a drawdown no deeper, and a better return in at least two of
the three stress periods.

\begin{table}[!htb]
\centering\small
\caption{Machine-learning variants, \res{sample_start}--\res{sample_end}, net of \res{cost_bps} per unit of turnover}\label{tab:ml}
\begin{tabularx}{\textwidth}{@{}Xrrrr@{}}
\toprule
\rowcolor{fmdark}\textcolor{white}{\textbf{Same months, net of costs}} & \textcolor{white}{\textbf{A: Trend}} & \textcolor{white}{\textbf{B: Vol.\ rule}} & \textcolor{white}{\textbf{C: HMM}} & \textcolor{white}{\textbf{D: Trend+HMM}} \\
CAGR (\% p.a.) & \resval{ml_a_cagr}{tab_ml}{4.6} & \resval{ml_b_cagr}{tab_ml}{5.3} & \resval{ml_c_cagr}{tab_ml}{4.9} & \resval{ml_d_cagr}{tab_ml}{4.2} \\
Volatility (\% p.a.) & \resval{ml_a_ann_vol}{tab_ml}{4.7} & \resval{ml_b_ann_vol}{tab_ml}{6.4} & \resval{ml_c_ann_vol}{tab_ml}{6.0} & \resval{ml_d_ann_vol}{tab_ml}{4.5} \\
Sharpe ratio & \resval{ml_a_sharpe}{tab_ml}{0.66} & \resval{ml_b_sharpe}{tab_ml}{0.61} & \resval{ml_c_sharpe}{tab_ml}{0.58} & \resval{ml_d_sharpe}{tab_ml}{0.61} \\
Maximum drawdown (\%) & \resval{ml_a_max_dd}{tab_ml}{$-$5.7} & \resval{ml_b_max_dd}{tab_ml}{$-$15.7} & \resval{ml_c_max_dd}{tab_ml}{$-$12.2} & \resval{ml_d_max_dd}{tab_ml}{$-$5.3} \\
Beta vs S\&P 500 & \resval{tab_ml}{tab_ml}{0.10} & \resval{tab_ml}{tab_ml}{0.22} & \resval{tab_ml}{tab_ml}{0.19} & \resval{tab_ml}{tab_ml}{0.07} \\
Turnover (\% p.a.) & \resval{ml_a_turnover}{tab_ml}{215.7} & \resval{ml_b_turnover}{tab_ml}{77.8} & \resval{ml_c_turnover}{tab_ml}{85.2} & \resval{ml_d_turnover}{tab_ml}{226.5} \\
Average T-bill weight (\%) & \resval{tab_ml}{tab_ml}{31.6} & \resval{tab_ml}{tab_ml}{3.6} & \resval{tab_ml}{tab_ml}{6.6} & \resval{tab_ml}{tab_ml}{34.5} \\
Stress periods better than A & -- & \resval{ml_b_stresswins}{tab_ml}{0 of 3} & \resval{ml_c_stresswins}{tab_ml}{0 of 3} & \resval{ml_d_stresswins}{tab_ml}{2 of 3} \\
Promotion gates passed & -- & \resval{ml_b_verdict}{tab_ml}{no} & \resval{ml_c_verdict}{tab_ml}{no} & \resval{ml_d_verdict}{tab_ml}{no} \\
\bottomrule
\end{tabularx}
\resmark{tab_ml}
\end{table}

\textbf{Result.} The HMM flags \res{ml_recall_hmm} of the \res{ml_n_stress} stress months (S\&P 500 down more than
\res{ml_stress_threshold} in the next month), more than the trend signal (\res{ml_recall_trend}). But only \res{ml_precision_hmm} of its
alarms are followed by such a fall, and it is risk-off in \res{ml_riskoff_hmm} of all months. Replacing the trend filter
(C) lowers the Sharpe ratio to \res{ml_c_sharpe} and deepens the drawdown to \res{ml_c_max_dd}, because the trend filter
also exits Treasuries when they fall. As an overlay (D) the HMM trims the drawdown to \res{ml_d_max_dd} but costs return
through extra cash and turnover (Sharpe \res{ml_d_sharpe}). No variant passes the gates (C \res{ml_c_verdict},
D \res{ml_d_verdict}); the result is the same for every random start and similar with three states
(\res{ml_3state_sharpe}). For our client this is a useful negative result: better crisis \emph{prediction} is not
automatically a better \emph{portfolio}, so the prospectus strategy stays rule-based.
